% English translation of checked_pages/esq_119_054.tex; source: cote 119, site 54, PDF page 55, typescript page 52.
% AI-translated by GPT-6 Astra (OpenAI, ChatGPT Pro mode).
\scanpage{119/54}
\mspage{52}

(\textsuperscript{1}) The expression ``out of reach'' here (and again further on for
an entirely different question), which I let pass despite
a reluctance, now strikes me as decidedly hasty and
unfounded. I have already been able to observe on other occasions that when
oracles (here myself !) declare with a knowing
(or doubtful) air that such-and-such a problem is ``out of reach'', this is
at bottom an entirely subjective assertion. It means
simply, apart from the fact that the problem is supposed not to have been
solved yet, that the speaker is short of ideas on the
question, or, probably more precisely, that he faces
it without feeling or enthusiasm, that it ``does nothing for him'', and
that he has no desire to do anything with it --- which
is often reason enough to want to discourage others
from doing so. This did not prevent the beautiful and lamented conjectures
of Mordell, Tate, and Chafarévitch, like Mr. de la Palisse,\tnote{TN-054-01}{The allusion evokes the proverbially obvious observation of a man still alive moments before dying. Here the joke concerns conjectures still declared out of reach just before they succumb. On the French allusion, see the entry ``truisme'' in the \href{https://oqlf.gouv.qc.ca/ressources/bibliotheque/dictionnaires/VocabulairesPDF/terminologie_figuresdestyle.pdf}{OQLF vocabulary of figures of speech}.}
from still being reputed ``out of reach'' at the very moment
of succumbing, the poor things ! --- Moreover, already in the days
following the writing of the present report, which brought me back
into contact with questions from which I had somewhat distanced myself
over the past year, I became aware of a new remarkable
property of the outer action of an absolute Galois
group on the fundamental group of an algebraic curve, which had
escaped me until now\frapparatus{V-FR-054-001} and which probably constitutes, at the very least,
a new step forward toward formulating an algebraic
characterization of $\mathrm{Gal}(\bar{\mathbb{Q}}/\mathbb{Q})$. This, together with the ``fundamental
conjecture'' (mentioned in par. 3 below), now appears as
the principal open question for the foundations of an ``anabelian
algebraic geometry'', which for several years has represented
(by far) my strongest focus of interest in mathematics.

(\textsuperscript{2}) I can nevertheless make an exception for another ``fact'', from the time
when, at around twelve years of age, I was interned in the concentration
camp of Rieucros (near Mende). It was there that I learned, from a
internee, Maria, who voluntarily gave me private lessons,

\margin{54'}
